\section{Ottimizzazione tramite Algoritmo Genetico}
\label{sec:ga}

\subsection{Motivazione}

Il raffinamento LM (Sezione~\ref{sec:distorsione}) opera in uno spazio
continuo e non può selezionare quali immagini includere né scegliere i
flag del solver, perché queste sono variabili discrete. Il problema
combinatorio --- quale sottoinsieme di \(m\) immagini e quale
combinazione di \(f\) flag minimizza \(\epsilon_{\text{RMSE}}\) --- ha
spazio di ricerca di dimensione \(2^m\cdot2^f\): per \(m=30\) e \(f=6\)
supera \(10^{10}\) configurazioni. Gli Algoritmi Genetici sono la
scelta naturale per esplorare questo spazio senza richiedere la
differenziabilità della funzione obiettivo~\cite{Marinescu2021}.

\subsection{Codifica del cromosoma}

Un cromosoma è una coppia:
\[
  c = \bigl(
    \underbrace{b_1,\dots,b_m}_{\text{maschera immagini}},\;
    \underbrace{f_1,\dots,f_k}_{\text{flag OpenCV}}
  \bigr),
  \quad b_i,f_j\in\{0,1\},
\]
con il vincolo \(\sum_i b_i \geq m_{\min}=5\) per garantire che il
solver abbia abbastanza viste.

\subsection{Funzione di fitness}

Per ogni cromosoma \(c\):
\begin{enumerate}[label=\arabic*.]
  \item Si selezionano le coppie
        \(\{(\homo{X}_i^{(j)},\homo{x}_i^{(j)})\}\) per gli indici
        attivi in \(c\).
  \item Si invoca il solver di calibrazione con i flag codificati in
        \(c\), ottenendo \(\mat{K}^c\), \(\mathbf{k}_r^c\) e i
        vettori di posa.
  \item Si calcola \(\epsilon_{\text{RMSE}}^c\) secondo
        \eqref{eq:reproj} per ogni vista attiva.
\end{enumerate}
La fitness è \(\epsilon_{\text{RMSE}}^c\) (da minimizzare).
Le valutazioni sono eseguite in parallelo su thread separati.

\subsection{Operatori genetici}

\paragraph{Crossover uniforme.}
Dati due genitori \(c_1,c_2\), ogni bit del figlio è scelto da
\(c_1\) o \(c_2\) con probabilità \(0.5\):
\[
  o_i =
  \begin{cases}
    (c_1)_i & \text{se } u_i < 0.5, \\
    (c_2)_i & \text{altrimenti,}
  \end{cases}
  \quad u_i\sim\mathcal{U}(0,1).
\]

\paragraph{Mutazione.}
Ogni bit viene invertito con probabilità \(p_m=0.25\), con il vincolo
\(\sum_i b_i\geq m_{\min}\).

\paragraph{Catastrofe (tipo 1).}
Se la diversità della popolazione scende sotto la soglia
\(\tau_\text{cat}=0.01\), una frazione \(p_\text{cat}=0.50\) viene
sostituita con cromosomi nuovi, prevenendo la convergenza prematura.

\paragraph{Catastrofe (tipo 2).}
Se la porzione inferiore della popolazione ha fitness \(=\infty\)
(solver divergente), la parte peggiore viene reinizializzata.

\subsection{Schema evolutivo}

\begin{figure}[H]
\centering
\begin{tikzpicture}[>=Stealth,
  node distance=0.85cm and 0cm,
  box/.style={draw, rounded corners=4pt, fill=yellow!10,
              minimum width=4.2cm, minimum height=0.8cm,
              align=center, font=\small},
  gbox/.style={draw, rounded corners=4pt, fill=green!12,
               minimum width=4.2cm, minimum height=0.8cm,
               align=center, font=\small},
  arr/.style={->, thick}]

  \node[box]  (init) {Inizializzazione pool (\(P=50\))};
  \node[box,  below=of init]  (eval) {Valutazione fitness (parallela)};
  \node[box,  below=of eval]  (sel)  {Selezione per ranking};
  \node[box,  below=of sel]   (xo)   {Crossover (\(p_c=0.75\))};
  \node[box,  below=of xo]    (mut)  {Mutazione (\(p_m=0.25\))};
  \node[box,  below=of mut]   (cat)  {Controllo diversità
                                       (catastrofe se necessario)};
  \node[box,  below=of cat]   (eden) {Archivia miglior cromosoma
                                       nell'Eden};
  \node[gbox, below=of eden]  (out)  {\(\mat{K}^*\),
                                       \(\mathbf{k}_r^*\),
                                       indici ottimali};

  \draw[arr] (init) -- (eval);
  \draw[arr] (eval) -- (sel);
  \draw[arr] (sel)  -- (xo);
  \draw[arr] (xo)   -- (mut);
  \draw[arr] (mut)  -- (cat);
  \draw[arr] (cat)  -- (eden);
  \draw[arr] (eden) -- (out)
        node[midway, right, font=\scriptsize] {epoch \(=\) max};

  % Ciclo
  \draw[arr] (eden.east) -- ++(1.2,0)
             -- ++(0,{0.85*5+0.25*4})
             -- (eval.east);
\end{tikzpicture}
\caption{Schema del ciclo evolutivo: il blocco Eden conserva il miglior
         cromosoma di ogni generazione; al termine si seleziona il
         miglior Eden.}
\label{fig:ga_schema}
\end{figure}